\section{Conclusion}
% In this paper, 
We introduce a novel metric called \ours\ for evaluating free-form and open-domain answers generated by large vision-language models. Compared to previous metrics, \ours\ offers a finer level of granularity, interpretability, and closer alignment with human judgments. Our quantitative analysis demonstrates that current LVLMs are prone to visual hallucination problems. We also find that the answer length and number of objects could affect the faithfulness of LVLMs. In addition, the faithfulness performance of LVLMs on different types of atomic facts varies. 
% In addition, our metric can also evaluate the traditional caption task. 
% Developing models that can obtain a high faithfulness performance across various types of attributes is a promising future research direction \xd{?}.
We expect that \ours\ will be of great value for evaluating forthcoming advanced LVLMs. 
% It's worth noting that, at present, \ours\ relies on ChatGPT for atomic facts generation, which can be expensive. Therefore, in the future, researchers can implement this metric using open-source models to make it more accessible and widely applicable. 

% It's worth noting that, at present, \ours relies on ChatGPT, which can be computationally expensive. However, in the future, we plan to implement this metric using open-source models to make it more accessible and widely applicable.
% \clearpage